\documentclass[12pt]{article}
\usepackage{amsmath, amssymb, graphicx, geometry, caption}
\geometry{margin=1in}
\setlength{\parindent}{0pt}


\title{Decrypting a Vigen\`ere Cipher with Known Key Length}
\author{Anurag Gupta}
\date{31 August 2025}

\begin{document}
\maketitle

\section{Introduction}
The goal was to decrypt an encrypted text that had been 
encoded with the \textbf{Vigen\`ere cipher}. 
The ciphertext consisted of lowercase alphabetic characters, 
and the \textbf{key length }was known in advance to be \textbf{5}. 
Using this information, I applied \textbf{frequency analysis} 
to guess the shifts of each key character, 
constructed the decryption key, and successfully recovered 
the text.

\section{Methodology}

\subsection{Reading Input Files}
The encrypted text was stored in a txt file and loaded into the script.

\subsection{Cipher Properties}
\begin{itemize}
    \item Cipher: Vigen\`ere (polyalphabetic substitution).
    \item Alphabet: lowercase English letters (a--z).
    \item Key length: 5.
\end{itemize}

\subsection{Frequency Analysis}
The ciphertext was divided into 5 sub-sequences by taking every 5th character, corresponding to each key position.

For each subsequence:
\begin{enumerate}
    \item Count letter frequencies using Python's \texttt{Counter}.
    \item Plot bar charts of letters vs. counts using \texttt{matplotlib}.
    \item Compare the distribution with known English frequencies (``e'' is most common, followed by ``t'', ``a''...).
\end{enumerate}

From this, I made educated guesses about which ciphertext peak corresponded to which plaintext letter.

\subsection{Key Guessing}
Suppose, in one subsequence, the most frequent letter was ``c''. If we assume it corresponds to ``e'' (most frequent in English), then the shift is:

\[
k = (\texttt{ord(c)} - \texttt{ord(e)}) \bmod 26 = 24 \rightarrow y
\]

This was repeated for all 5 positions.  

\begin{center}
\begin{tabular}{|c|c|c|c|c|}
\hline
Position & Assumed Plaintext & Cipher Peak & Shift & Letter \\
\hline
0 & e & c & (c-e) mod 26 = 24 & y\\
1 & e & e & (e-e) mod 26 = 0 & a\\
2 & e & q & (q-e) mod 26 = 12 & m\\
3 & t & d & (d-t) mod 26 = 10 & k\\
4 & e & b & (b-e) mod 26 = 23 & x\\
\hline
\end{tabular}
\end{center}

This produced a letter key which is \textbf{yamkx}

\subsection{Decryption}
Once the key was determined, the decryption process was straightforward:

For each ciphertext character $c_i$ at position $i$ we get a decrypted charachter $d_i$:
\[
d_i = \Big( ord(c_i) - ord(k_{i \bmod 5})  \Big) \bmod 26 + ord(a)
\]

This was implemented in the \texttt{decrypt\_text} function.

\subsection{Output}
Two files were written:
\begin{itemize}
    \item Decrypted plaintext (\texttt{decrypt.txt})
    \item Guessed key (\texttt{key\_guess.txt})
\end{itemize}

\section{Results}
\begin{itemize}
    \item The frequency plots revealed the relative shifts of each cipher.
    \item By aligning ciphertext peaks with expected plaintext peaks, and guessing if the
      decrypted english text makes sense I was able to guess a 5-character key.
    \item The ciphertext was successfully decrypted using the key guessed from the relative shifts.
\end{itemize}

\section{Frequency Analysis}

This is the freequency plots for the encrypted text (${c_i \bmod 5}$).

\includegraphics[width=0.5\textwidth]{/mnt/Personal/Projects/Micro_Codes/Python/Crypto/Assignment_1/plots/p1.png}
\includegraphics[width=0.5\textwidth]{/mnt/Personal/Projects/Micro_Codes/Python/Crypto/Assignment_1/plots/p2.png}
\includegraphics[width=0.5\textwidth]{/mnt/Personal/Projects/Micro_Codes/Python/Crypto/Assignment_1/plots/p3.png}
\includegraphics[width=0.5\textwidth]{/mnt/Personal/Projects/Micro_Codes/Python/Crypto/Assignment_1/plots/p4.png}
\includegraphics[width=0.5\textwidth]{/mnt/Personal/Projects/Micro_Codes/Python/Crypto/Assignment_1/plots/p5.png}

\section{Output Analysis}

Thisis the alphabet frequency plot of my decrypted message.

\includegraphics[width=0.75\textwidth]{/mnt/Personal/Projects/Micro_Codes/Python/Crypto/Assignment_1/plots/final_plot.png}


This is the plot of average letter frequncies in english from wikipedia.

\includegraphics[width=0.8\textwidth]{/mnt/Personal/Projects/Micro_Codes/Python/Crypto/Assignment_1/plots/normal.png}

As we can see that my plot closely represents the expected plot from wikipedia which shows 
that my output is a proper decrypted english message.

\end{document}
